\section{15.3 멀티에이전트 RL 개관 (선택)}
\begin{frame}{에이전트가 하나가 아니라면}
  \begin{itemize}
    \item 지금까지는 \textbf{에이전트가 하나뿐인} 환경을 다뤘다
    \item \kb{멀티에이전트 강화학습(Multi-Agent RL, MARL)}:
      여러 에이전트가 \textbf{같은 환경에서 동시에} 행동
      (협력: 팀 스포츠·물류 로봇 협업, 경쟁: 바둑·대전 게임 AI)
    \item 근본적 차이: 환경이 더 이상 \textbf{고정된 규칙만} 따르지
      않음 --- 다른 에이전트도 학습하며 계속 전략을 바꿈
  \end{itemize}
  \vskip0.6em
  \begin{block}{움직이는 표적(non-stationary target)}
    ``\textbf{내 입장에서 본 환경}''이 학습 도중에 \textbf{계속
    움직이는 표적} --- Week 3.1의 ``환경은 고정된
    $P(s'|s,a)$를 따른다''는 \textbf{마르코프 성질의 전제}
    자체가, 다른 에이전트를 환경의 일부로 보는 순간 흔들림.
  \end{block}
  \vskip0.5em
  \begin{itemize}
    \item 단일 에이전트 MDP: 기대리턴은 \textbf{자기 정책만의
      함수} $J(\pi) = \mathbb{E}_{\tau \sim \pi}\left[\sum_t
      \gamma^t r_t\right]$
    \item 에이전트가 둘이면: $J_1(\pi_1, \pi_2)$,
      $J_2(\pi_1, \pi_2)$ --- \textbf{자기 정책과 나머지 모든
      에이전트의 정책}의 함수
  \end{itemize}
\end{frame}
